% Book-local P06.1 relative positive-weighted-cone result, 10 September 2026.
% It records a direct local calculation; no priority claim is made.
\begin{proposition}[Pluriharmonic descent along a relative positive-weighted cone]
\label{prop:p061-relative-positive-weighted-cone-descent}
Let \(K\) be a complete non-Archimedean field with nontrivial absolute
value, let \(B\) be a strict \(K\)-affinoid algebra, and put
\(Y:=\mathcal M(B)\).  Fix positive integers \(w_1,\ldots,w_m\), and let
\[
 F_1,\ldots,F_r\in B[T_1,\ldots,T_m]
\]
be homogeneous of strictly positive weighted degree, where \(B\) has weight
zero and \(T_i\) has weight \(w_i\).  Write
\[
 X:=\mathcal M\bigl(B\langle T_1,\ldots,T_m\rangle/(F_1,\ldots,F_r)\bigr),
 \qquad p:X\longrightarrow Y,
\]
and let \(s:Y\to X\) be the zero section.  For every \(y\in Y\), every
ordinary open neighbourhood \(\Omega\) of \(s(y)\), and every classically
pluriharmonic function \(h:\Omega\to\mathbf R\), there are a strict
affinoid ordinary neighbourhood \(V\) of \(y\) in \(Y\), an integer
\(n\geq1\), and \(\varpi\in K\) with \(0<|\varpi|<1\), such that the
strict affinoid domain
\[
 U_n(V):=\left\{x\in p^{-1}(V):
 |T_i(x)|\leq |\varpi|^{n w_i}\quad(1\leq i\leq m)\right\}
 \subset\Omega
\]
has \(s(y)\) in its ordinary interior and satisfies
\begin{equation}
\label{eq:p061-relative-weighted-cone-factorization}
 h|_{U_n(V)}=p^*g,
 \qquad g:=s^*h|_V\in\operatorname{PH}(V).
\end{equation}
If \(Y\) is quasi-smooth, then \(h\) is locally \(\mathbf R\)-PL and
Gubler--Jell--Rabinoff harmonic at \(s(y)\).
\end{proposition}
